\section{Fourfold examples and the role of monodromy}\label{sec:examples}

\subsection{The isolated and interacting plane contractions}
The compact fourfolds in \cite{wuebben17} provide two further
applications of the general theorem. We verify its local
simple-zero hypothesis explicitly and record the additional
global metric conclusion; their construction and detailed cycle
geometry remain in that paper.

For the isolated example, the lattice is
\(\NS(X)=\Z H\oplus\Z D\) with Gram matrix
\(\operatorname{diag}(2,-22)\), and
\[
 a=(1,0,1),\quad b=(2,D,-5),\quad v=a+b.
\]
The compact chamber spaces have a single \(\PP^2\) fibre.
The complete factor list and no-charge theorem
\cite[Section~3]{wuebben17} verify (G1)--(G2).
The full deep representation is
\(\C^3\oplus(\C^3)^*\) under \(U(1)\), with complex equation
\(y(x)=0\). The point \(x=e_1,y=e_2^*\) is module-simple,
as both arrow directions between the two one-dimensional
vertices are nonzero. At stable points the effective group
is trivial. The convergent charts and classifying diagram
verify (G3).

For the interacting example, the lattice matrix and vectors are
\[
 \begin{pmatrix}2&0&0\\0&-6&-4\\0&-4&-8\end{pmatrix},\qquad
 a_0=(1,0,1),\quad a_1=(1,D_1,-2),\quad a_2=(1,D_2,-3),
 \quad v=a_0+a_1+a_2 .
\]
For \([\omega]=H+sD_1+tD_2\), the levels are
\[
 \zeta(s,t)=\frac{2\pi}{3}
 \begin{pmatrix}-10&-12\\8&0\\2&12\end{pmatrix}
 \binom{s}{t}.
\]
All six chamber spaces, their common normal target, complete
factor and charge lists, and whole fibres are constructed in
\cite[Sections~4--5]{wuebben17}. They verify (G1)--(G2).
At the deepest point the full effective group is
\(U(1)^3/U(1)\), with paired arrow dimensions \(1,2,1\).
In the coordinates of that paper,
\[
 zz'-\varphi(u)=ww'+\varphi(u)=0
\]
has the solution \(z=z'=w=1,\ w'=-1,\ u=\varphi=(1,0)\).
Its directed support is strongly connected, which for
one-dimensional vertices forbids a proper submodule.
Along the moving stratum the full representation is
\(\C^2\) times the two-arrow \(U(1)\) model, which has
the simple point \(e_1,e_2^*\); stable centers are smooth.
Together with the convergent classifying diagrams this verifies
(G3) on every stratum.

Theorem~\ref{thm:quantitative} consequently gives graph distortion
tending to zero in both examples, with the variational bound
\eqref{eq:global-distortion}. For the interacting family this
holds through all chamber directions and arbitrary ratios of
positive scales. It supplements the qualitative global
convergence of \cite{wuebben17} by a proof from the finite
quotient distances.

Their integral formulas are instances of
Theorem~\ref{thm:attachment}, with the explicit normalized
index restrictions already proved in
\cite[Proposition~6.1]{wuebben17}.
For the plane, \(B=(3)\) and
\[
 \kappa(c_2(I_z)|_V)=\binom{m_z}{2},\qquad
 \kappa(c_1(I_z)c_1(I_{z'})|_V)=m_zm_{z'}\pmod3.
\]
Here \(m_z\) is the Euler pairing with the extension subobject.
In the two-plane chamber, \(E=P_0\cup P_2\) and
\[
 B=\begin{pmatrix}3&1\\1&3\end{pmatrix},\qquad
 \coker B\simeq\Z/8,\qquad (n_0,n_2)\longmapsto n_0-3n_2 .
\]
With \(m_{i,z}=\chi(a_i,z)\), the boundary map is
\[
 \begin{aligned}
 \kappa(c_2(I_z)|_V)
   &=\binom{m_{0,z}}2-3\binom{m_{2,z}}2,\\
 \kappa(c_1(I_z)c_1(I_{z'})|_V)
   &=m_{0,z}m_{0,z'}-3m_{2,z}m_{2,z'}\pmod8 .
 \end{aligned}
\]
These formulas use the marked complement ring and respect all
its relations. The mod-eight statement concerns this
two-plane chamber. The other four configurations are not
assigned its global formula.

The proper moving-line family specializes as \(f=c_0+c_2\),
with both multiplicities one. For levels
\(\zeta_0=a>0,\zeta_2=b>0\), the exact areas are
\(\pi a,\pi b,\pi(a+b)\); the plane volumes are
\((\pi a)^2/2,(\pi b)^2/2\).
The ambient collapse and complete separated-scale limits are
those of \cite[Section~8]{wuebben17}, with unchanged metric
normalization. Thus the general global comparison and the
explicit attachment refer to those same cycles.

\subsection{The Schubert involution}
The following topological example explains why the geometric
transition information in Section~\ref{subsec:retained} cannot
be replaced by local fibre rings and divisor classes.

Choose a nondegenerate symmetric complex bilinear form on
\(\C^4\). Orthogonal complementation defines a holomorphic
involution \(\delta(P)=P^\perp\) of \(G=\Gr(2,4)\).
For its tautological subbundle \(\mathcal S\) and quotient
\(\mathcal U\), one has
\(\delta^*\mathcal S\simeq\mathcal U^*\). Consequently
\[
 \delta^*h=h,\qquad \delta^*p=q,\qquad \delta^*q=p .
\]
Form the compact smooth mapping tori \(T_1\) and \(T_\delta\)
over \(S^1\). Both have the same local product geometry,
integral fibre ring, and distinguished divisor class \(h\).
Their degree-four cohomology is different:
\begin{equation}\label{eq:monodromy}
 H^4(T_1,\Z)=\Z^2,\qquad H^4(T_\delta,\Z)=\Z .
\end{equation}
Indeed \(H^3(G,\Z)=0\), so the Wang sequence identifies
degree-four cohomology with the invariant subgroup of
\(H^4(G,\Z)=\Z p\oplus\Z q\). For the identity it is all of
that group; for \(\delta\) it is \(\Z(p+q)\).

This is a topological necessity example, not a claim that
the involution occurs in the K3 family above. It shows that
forgetting monodromy can change even a Betti number.
The use of all multiple intersections has a further purpose:
cochains retain their restriction maps and higher compatibility,
which fibre cohomology alone generally discards.
Theorem~\ref{thm:attachment} gives a smaller calculation only
under its explicit vanishing, intersection and generator
hypotheses.

The results concern compact K3 bundle-moduli contractions
without bubbling. Extending them to positive-charge strata
requires new local geometric and metric inputs. The
reconstruction of underlying degenerating manifolds and
unrestricted Hermitian--Yang--Mills dynamics are likewise
separate questions.
